% 08_biophysics.tex
% Ch.8: Protein backbone impedance

\chapter{Biophysics: Protein Folding}
\label{ch:biophysics}

\section{Protein Backbone: From Proton Radius to Folding}
\label{sec:protein_bridge}

The protein $S_{11}$ folding engine uses the C$\alpha$--C$\alpha$
virtual bond distance $d_0 \approx 3.80$~\AA\ as its fundamental
length scale.  This traces back to the AVE axioms via:
\begin{resultbox}{Proton-to-Backbone Length Scale}
\begin{equation}
    d_p = \frac{4\hbar}{m_p c} \approx 0.841\;\text{fm}
    \;\xrightarrow{\text{nuclear bond}}
    r_{\text{cov}}
    \;\xrightarrow{\text{backbone}}
    d_0 \approx 3.80\;\text{\AA}
\end{equation}
\end{resultbox}
The ratio $d_0 / a_0 \approx 7.18$ (where $a_0 = \ell_{\text{node}}
/\alpha$ is the Bohr radius) suggests a geometric origin in the
peptide repeat unit.  The amino acid impedance table $Z_{\text{topo}}$
assigns complex impedances (R + jX) to each of the 20 residues,
with $Q_{\text{backbone}} \approx 7$ suppressing reactive coupling.

\begin{figure}[H]
    \centering
    \includegraphics[width=0.95\textwidth]{protein_backbone_impedance.pdf}
    \caption{Protein backbone impedance from AVE axioms.
    \textbf{Left:} Full derivation chain from $\ell_{node}$ to
    the C$\alpha$--C$\alpha$ peptide bond length $d_0 = 3.80$~\AA.
    \textbf{Centre:} Amino acid impedance spectrum --- hydrophobic
    (gold), positive/inductive (blue), negative/capacitive (red).
    \textbf{Right:} Local inter-residue reflection coefficient
    $\Gamma$ along a sample peptide sequence; high $|\Gamma|$
    indicates folding-driving impedance mismatch.}
    \label{fig:protein_backbone}
\end{figure}
